\scanpage{123}
\mspage{110.}
\inserted{(N\uline{B}.} \deleted{\unsure{Par Zorn.}} $\unsure{\mathcal H|H}$ n'est pas vide, car $(\varnothing,h_\varnothing)$ est
dans $\mathcal H$. -- mais c'est contenu dans l'\deleted{\unsure{énoncé}}
énoncé précédent que tte partie \unsure{tot} ordonnée
de $\mathcal H$ y a une borne sup\ldots) Par Zorn
il existe un élément maximal $(U,h_U)$.
À prouver que $U=I$, et on aura terminé.
\inserted{On \unsure{va voir|verra} que si $(U,h_U)$ tel que $U\ne I$, ce n'est \uline{pas} maximal.}
Soit $\deleted{\unsure{B}}F=I\setminus U\ne\varnothing$ le fermé complémentaire.
Par l'hypothèse artinienne sur $I$,
$F$ a un élément minimal $i_0$. \deleted{On}
Soit
\[
 U'=U\cup\overbrace{I_{/\unsure{j_0|i_0}}}^{V}=U\cup\{i_0\},
\]
\deleted{on} \deleted{\unsure{doit}}\inserted{\unsure{si}} prolonger $h_U$ en $h_{U'}$ sur $U'$,
on aura terminé. Pour ceci, soit
\deleted{on est ramené à faire la}
$V_0=V\cap U=V\setminus\{j_0\}$, on est ramené
à prolonger la restriction \inserted{$h_{V_0}$} de $h_U$ à $V_0$,
de $V_0$ à $V$. Au total, \unsure{OPS}
$I=I_{/i_0}$ muni d'un élément final $i_0$,
$U=I\setminus\{i_0\}$, et montrer qu'un
$(U,h_U)\in\mathcal H$ se prolonge à $I$. Soient
$A',B',X',Y'$ les restrictions de $A,B,X,Y$
à $U$ \deleted{\unsure{et}} (qu'on pourra aussi noter $I'$).
La donnée de $A$ équivaut à celle
de $A'$, celle de $A_0=A(i_0)\in\operatorname{Ob} M$, et
d'une flèche de
\[
 A'_{I'}\overset{\mathrm{def}}{=}
 \underset{i'\in I'}{\underrightarrow{\lim}}\,A'(i')\xrightarrow{\unsure{\lambda_A|\partial_A}}
 A_0\bigl(\simeq\underset{\unsure{\nearrow}}{A_I}\bigr),
\]
% archive sheet number: 122

